\documentclass[12pt]{article}
\usepackage[utf8]{inputenc}
\usepackage{graphicx}
\usepackage[numbers]{natbib}
\usepackage{amsmath}
\usepackage{amssymb}
\usepackage{mathrsfs}
\usepackage{hyperref}
\usepackage{natbib}
\usepackage{listings}
\usepackage{enumitem}
\usepackage{pgfplots}
\pgfplotsset{compat=1.18}
\usepackage{multirow}
\usepackage{booktabs}
\usepackage{algorithm}
\usepackage{algpseudocode}
\usepackage{tikz}
\usepackage{tabularx}
\usepackage{verbatim}
\usetikzlibrary{positioning,fit,backgrounds,arrows.meta,calc}

% Load packages that style.sty uses but doesn't load
\usepackage{geometry}
\usepackage{xcolor}
\usepackage{fancyhdr}
\usepackage{lastpage}

% Load style file (uses: geometry, xcolor, titlesec, titling, fancyhdr commands)
\usepackage{style}

% --- explicit layers so backgrounds never cover content ---
\pgfdeclarelayer{bg}
\pgfdeclarelayer{fg}
\pgfsetlayers{bg,main,fg}

% Define firstpage style (style.sty references it but Overleaf needs it defined explicitly)
\fancypagestyle{firstpage}{
  \fancyhf{}
  \fancyfoot[R]{\thepage/\pageref*{LastPage}}
}

% Math operators
\DeclareMathOperator*{\argmin}{arg\,min}
\DeclareMathOperator*{\argmax}{arg\,max}

% Code style
\lstset{
    basicstyle=\ttfamily\small,
    keywordstyle=\color{blue},
    commentstyle=\color{green!60!black},
    stringstyle=\color{red},
    numbers=left,
    numberstyle=\tiny,
    frame=single,
    breaklines=true,
    backgroundcolor=\color{gray!10}
}

% Slide-optimized settings
\titleformat{\section}
  {\Large\bfseries\color{TAMUmaroon}}
  {}{0pt}{}

\setlength{\parskip}{0.5em}
\setlength{\itemsep}{0.2em}
\pagestyle{fancy}

\title{\Large Research Update}
\author{Gaoming Lin \\
\vspace{10pt} Advisor: Dr. Byung-Jun Yoon}
\date{01/30/2026}

\begin{document}

\maketitle
\thispagestyle{firstpage}

% =========================================================
% Parameter-Driven Literature Slides (Paper 1–5)
% =========================================================

\newpage
\section*{Paper 1: NREL 2024 — Probing Signal-Based Inertia \& Frequency Response Estimation \cite{peng2024nrel}}

\textbf{What the paper did}
\begin{itemize}
  \item Injects active-power probing signals through inverter-based resources (IBRs).
  \item Shapes probing signals in the frequency domain to isolate inertial dynamics.
  \item Estimates inertia and droop using swing-equation-based regression.
\end{itemize}

\textbf{Probing parameter guidance extracted from this paper}
\begin{itemize}
  \item \textbf{Amplitude ($A$):}
  \begin{itemize}
    \item Must exceed ambient frequency noise to ensure measurable ROCOF.
    \item Upper bounded by grid-security constraints (no protection triggering).
  \end{itemize}

  \item \textbf{Frequency band ($\omega$):}
  \begin{itemize}
    \item Probing energy concentrated in low-frequency range where inertia dominates.
    \item Avoids overlap with primary frequency control and inverter inner loops.
  \end{itemize}

  \item \textbf{Duration ($T$):}
  \begin{itemize}
    \item Long enough to observe inertial transient.
    \item Short relative to slower governor and secondary control responses.
  \end{itemize}
\end{itemize}

\textbf{Why this is critical for my MOCU-based sBOED}
\begin{itemize}
  \item Defines physically meaningful bounds on $(A, T, \omega)$.
  \item Shows probing should maximize \textbf{sensitivity of ROCOF to inertia}, not generic excitation.
  \item Directly defines the sBOED design variable:
  \[
  u = \{A,\ T,\ \omega\}.
  \]
\end{itemize}




% ---------------------------------------------------------

\newpage
\section*{Paper 2: HICSS 2024 — Real-Time Inertia Estimation Using Probing \& PMUs \cite{yi2024hicss}}

\textbf{What the paper did}
\begin{itemize}
  \item Uses repeated low-magnitude probing signals for online inertia estimation.
  \item Identifies inertia via PMU-measured frequency deviations.
  \item Emphasizes secure and non-disruptive probing.
\end{itemize}

\textbf{Probing parameter guidance extracted from this paper}
\begin{itemize}
  \item \textbf{Amplitude ($A$):}
  \begin{itemize}
    \item Explicitly limited to small active-power deviations.
    \item Ensures probing does not violate frequency limits or stability margins.
  \end{itemize}

  \item \textbf{Observation window ($T$):}
  \begin{itemize}
    \item Chosen to isolate inertial response before governor action.
    \item Too long a window contaminates inertia estimates with control effects.
  \end{itemize}

  \item \textbf{Repetition strategy:}
  \begin{itemize}
    \item Multiple probes allow belief refinement without increasing $A$.
  \end{itemize}
\end{itemize}

\textbf{Why this is critical for my MOCU-based sBOED}
\begin{itemize}
  \item Provides realistic numerical ranges for $A$ usable as hard constraints.
  \item Motivates \textbf{sequential probing}: adapt $(A,T)$ as uncertainty shrinks.
  \item Reinforces probing as a constrained decision problem, not free excitation.
\end{itemize}




% ---------------------------------------------------------

\newpage
\section*{Paper 3: PMU + Swing Equation — Real-Time Inertia Estimation \& Validation (2025) \cite{pmu2025}}

\textbf{What the paper did}
\begin{itemize}
  \item Uses the swing equation as the latent dynamic model.
  \item Estimates inertia from PMU frequency and ROCOF.
  \item Validates inferred inertia against system benchmarks.
\end{itemize}

\textbf{Probing parameter implications (indirect but crucial)}
\begin{itemize}
  \item \textbf{Target observable:}
  \begin{itemize}
    \item ROCOF is most informative for inertia.
    \item Probing must induce sufficient acceleration, not just steady deviation.
  \end{itemize}

  \item \textbf{Timing ($T$):}
  \begin{itemize}
    \item Probing must excite the initial transient phase.
    \item Slow probing yields weak ROCOF and poor identifiability.
  \end{itemize}
\end{itemize}

\textbf{Why this is critical for my MOCU-based sBOED}
\begin{itemize}
  \item Defines what probing must excite to reduce inertia uncertainty.
  \item Guides utility design: experiments that fail to induce ROCOF have low MOCU value.
  \item Strengthens the observation model used for Bayesian belief updates.
\end{itemize}




% ---------------------------------------------------------

\newpage
\section*{Paper 4: Du et al. 2022 — Approximations for Optimal Experimental Design \cite{du2022}}

\textbf{What the paper did}
\begin{itemize}
  \item Applies Fisher-information-based OED to power-system parameter estimation.
  \item Selects excitation inputs to maximize parameter identifiability.
\end{itemize}

\textbf{Probing parameter insights from classical OED}
\begin{itemize}
  \item Treats probing parameters as optimization variables under constraints.
  \item Encourages excitation directions that maximize parameter sensitivity.
\end{itemize}

\textbf{Why this is critical for my MOCU-based sBOED}
\begin{itemize}
  \item Provides a baseline: how $(A,T,\omega)$ would be chosen if estimation accuracy were the only goal.
  \item Highlights limitation: information-maximizing probes may not improve stability decisions.
  \item Motivates replacing FIM with \textbf{MOCU-based utility} for probing parameter selection.
\end{itemize}




% ---------------------------------------------------------

\newpage
\section*{Paper 5: IJEPES 2024 — Global Inertia Estimation via Active Perturbations \cite{brambilla2024}}

\textbf{What the paper did}
\begin{itemize}
  \item Injects probing tones using grid-forming converters.
  \item Estimates global inertia from frequency-domain response.
  \item Models the grid as a synchronized oscillator.
\end{itemize}

\textbf{Probing parameter guidance extracted from this paper}
\begin{itemize}
  \item \textbf{Frequency selection:}
  \begin{itemize}
    \item Probing tones placed in synchronized low-frequency modes.
    \item High-frequency excitation provides little inertia information.
  \end{itemize}

  \item \textbf{Signal structure:}
  \begin{itemize}
    \item Frequency-domain probing improves identifiability of inertial constants.
  \end{itemize}
\end{itemize}

\textbf{Why this is critical for my MOCU-based sBOED}
\begin{itemize}
  \item Supports selecting probing frequency content based on system-level objectives.
  \item Justifies global-objective-driven probing (e.g., ROCOF risk).
  \item Reinforces probing design as an objective-based decision problem.
\end{itemize}



\newpage
\bibliographystyle{ieeetr}
\bibliography{references}
\end{document}
